\begin{abstract}

Software engineers create programming languages for specific purposes, and
effective usage of the programming language requires effective editing support.
Announced in 2016, language server protocol aims to help software engineers 
obtain enriched support in their text editor (vim, neovim) or IDE (VSCode,
Eclipse IDE) of choice by decoupling the `language' and the `editing' parts.
Since modern IDEs and some text editors support language server protocol `out of
the box', it suffices for engineers to implement only one piece of software
called language server to gain some editing features. For some reasons, not all
software languages get language server support. One of these languages is
Jasmin. Jasmin is a cryptographic primitive description language and the primary 
users of this language are the cryptography researchers. In this work, I present
an implementation of a language server for the Jasmin programming language,
which supports diagnostics, goto definition, incremental parsing, symbol
highlighting, syntax highlighting, and symbol type on hover.

\end{abstract}